%%%%%%%%%%%%%%%%%
% Numerische Ergebnisse
%%%%%%%%%%%%%%%%%
\section{Results}
All experiments were conducted on an \textbf{AMD Ryzen 7 5700X} CPU (8 cores, 32\,MB L3 cache, measured STREAM triad bandwidth of $36.2\,\text{GB/s}$) and an \textbf{AMD Radeon RX 9070} GPU (measured STREAM triad bandwidth of $682.45\,\text{GB/s}$). 

Across both architectures, all three stencils ($S_5$, $S_9$, and $S_5^{\text{diag}}$) achieve nearly identical throughput, confirming that the Jacobi kernels are strictly \emph{memory-bound}. On the large $2^{24}$-point grid, the spatially blocked CPU implementation sustains $\approx 1500\,\text{MUpdates/s}$, effectively saturating the Ryzen 7 5700X's main memory bandwidth. Without spatial blocking, CPU performance degrades by up to $40\%$ (dropping to $\approx 891\,\text{MUpdates/s}$) once the row length reaches $N_x \ge 65\,536$. At this size, three consecutive grid rows exceed the L1/L2 cache capacity, leading to cache eviction and redundant DRAM accesses. Spatial blocking eliminates this layer-condition bottleneck entirely.

On the Radeon RX 9070, the pure compute kernel maintains a stable throughput of $\approx 13.2$ GUpdates/s across all grid aspect ratios---an $\approx 8.8\times$ speedup over the blocked CPU execution, driven by the GPU's $18.85\times$ higher memory bandwidth. While PCIe host-to-device data transfers noticeably reduce the effective throughput during short 50-sweep runs ($6.0\text{--}10.9\,\text{GUpdates/s}$), this transfer overhead becomes negligible ($< 2\%$) over $1000$ sweeps. On the smaller $512 \times 512$ grid ($2\,\text{MB}$ per double-precision array), the working set fits entirely within the Ryzen 7 5700X's $32\,\text{MB}$ L3 cache, boosting CPU throughput to $\approx 5.0\text{--}5.4\,\text{GUpdates/s}$, whereas the GPU reaches $\approx 10.5\,\text{GUpdates/s}$ due to kernel launch latencies.

Finally, comparing numerical convergence on the $512 \times 512$ grid after $1000$ sweeps highlights the efficiency of the compact 9-point stencil $S_9$: because its additional arithmetic operations are hidden by the memory-bandwidth limit, $S_9$ incurs no runtime penalty while reducing the $L_2$ residual by more than two orders of magnitude compared to $S_5$ ($0.099$ vs.\ $12.76$).

\subsection*{Performance and Convergence Analysis}

All experiments were conducted on an \textbf{AMD Ryzen~7~5700X} CPU (8 cores, $32\,\text{MB}$ L3 cache, measured STREAM triad bandwidth of $37.4\,\text{GB/s}$) and an \textbf{AMD Radeon~RX~9070} GPU ($8\,\text{MB}$ L2 cache, $64\,\text{MB}$ Infinity Cache, measured STREAM triad bandwidth of $684.28\,\text{GB/s}$).

On the large $2^{24}$-point grid ($128\,\text{MB}$ per double-precision array), all three stencils ($S_5$, $S_9$, and $S_5^{\text{diag}}$) achieve nearly identical throughput across both architectures, confirming that the kernels are strictly \emph{memory-bound} (Table~\ref{tab:aspect_ratio}). The spatially blocked CPU version sustains $\approx 1.50\text{--}1.54\,\text{GUpdates/s}$, saturating the $37.4\,\text{GB/s}$ main memory bandwidth at the theoretical minimum of $24\,\text{Bytes/update}$. Without spatial blocking, CPU performance drops by up to $41\%$ (to $\approx 0.89\,\text{GUpdates/s}$) for row lengths $N_x \ge 65\,536$, as three consecutive grid rows exceed the L1/L2 cache capacity and trigger cache evictions. Spatial blocking eliminates this bottleneck completely.

\begin{table}[htbp]
    \centering
    \caption{Throughput (in GUpdates/s) across varying grid aspect ratios ($N_x \times N_y = 2^{24}$) for the compact 9-point stencil $S_9$ over $50$ sweeps.}
    \label{tab:aspect_ratio}
    \small
    \begin{tabular}{rrcccc}
        \hline
        \textbf{$N_x$} & \textbf{$N_y$} & \textbf{CPU (unblocked)} & \textbf{CPU (blocked)} & \textbf{GPU (pure)} & \textbf{GPU (transfer)} \\
        \hline
        $256$         & $65\,536$ & $1.53$ & $1.54$ & $13.16$ & $11.21$ \\
        $4096$        & $4096$    & $1.51$ & $1.51$ & $13.14$ & $10.23$ \\
        $32\,768$     & $512$     & $1.32$ & $1.32$ & $13.13$ & $10.91$ \\
        $65\,536$     & $256$     & $1.10$ & $1.42$ & $13.12$ & $10.28$ \\
        $262\,144$    & $64$      & $0.89$ & $1.42$ & $13.22$ & $10.91$ \\
        $2\,097\,152$ & $8$       & $1.38$ & $1.42$ & $13.15$ & $9.63$  \\
        \hline
    \end{tabular}
\end{table}

On the Radeon~RX~9070, the pure compute kernel maintains a stable throughput of $\approx 13.1\text{--}13.4\,\text{GUpdates/s}$ across all aspect ratios on the $2^{24}$ grid---an $\approx 8.8\times$ speedup over the blocked CPU implementation. While PCIe host-to-device transfers reduce effective throughput during short 50-sweep runs ($9.6\text{--}11.4\,\text{GUpdates/s}$), this overhead drops to $< 3.5\%$ over $1000$ sweeps (Table~\ref{tab:convergence_perf}).

\begin{table}[htbp]
    \centering
    \caption{Throughput (in GUpdates/s) and final $L_2$ residual norm after $1000$ Jacobi sweeps on the Ryzen~7~5700X (blocked) and Radeon~RX~9070.}
    \label{tab:convergence_perf}
    \small
    \begin{tabular}{llcccc}
        \hline
        \textbf{Grid Size} & \textbf{Stencil} & \textbf{CPU} & \textbf{GPU (pure)} & \textbf{GPU (total)} & \textbf{$L_2$ Residual} \\
        \hline
        $512 \times 512$   & $S_5$               & $4.93$ & $10.47$ & $10.27$ & $12.764$ \\
                           & $S_9$               & $5.25$ & $10.48$ & $10.36$ & $0.099$  \\
                           & $S_5^{\text{diag}}$ & $5.36$ & $10.70$ & $10.56$ & $6.914$  \\
        \hline
        $2048 \times 2048$ & $S_5$               & $2.52$ & $13.41$ & $12.95$ & $57.767$ \\
                           & $S_9$               & $2.54$ & $16.56$ & $16.39$ & $0.199$  \\
                           & $S_5^{\text{diag}}$ & $2.52$ & $22.18$ & $21.88$ & $29.237$ \\
        \hline
    \end{tabular}
\end{table}

Comparing the $512 \times 512$ ($2\,\text{MB}$/array) and $2048 \times 2048$ ($32\,\text{MB}$/array) grids over $1000$ sweeps reveals distinct cache and latency regimes:
\begin{itemize}
    \item \textbf{CPU Cache Hierarchy:} On the $512 \times 512$ grid, the working set fits entirely inside the Ryzen~7~5700X's $32\,\text{MB}$ L3 cache, yielding $\approx 4.9\text{--}5.4\,\text{GUpdates/s}$. At $2048 \times 2048$, the arrays partially spill into main memory, settling at an intermediate rate of $\approx 2.52\,\text{GUpdates/s}$.
    \item \textbf{GPU Launch Latency vs.\ Infinity Cache:} Despite fitting within the GPU's $8\,\text{MB}$ L2 and $64\,\text{MB}$ L3 caches, the $512 \times 512$ grid is limited to $\approx 10.5\,\text{GUpdates/s}$ by kernel launch overhead ($\approx 25\,\mu\text{s}$ per sweep) and low compute-unit occupancy ($262\text{k}$ threads). At $2048 \times 2048$ ($4.19\text{M}$ threads), launch latencies are amortized and the working set still benefits from the $64\,\text{MB}$ Infinity Cache, pushing throughput up to $13.4\text{--}22.2\,\text{GUpdates/s}$ (up to an $8.8\times$ GPU speedup).
    \item \textbf{Numerical Convergence:} Across both grid resolutions, the compact 9-point stencil $S_9$ incurs no runtime penalty over $S_5$ while reducing the final $L_2$ residual after $1000$ sweeps by more than two orders of magnitude ($0.199$ vs.\ $57.767$ at $2048 \times 2048$).
\end{itemize}

